\section*{Abstract}



Die Natur und ihre Bedürfnisse sind für Menschen nicht immer auf den ersten Blick ersichtlich. Um die Kommunikation und Bewusstheit der Menschen ihrem Umfeld näher zu bringen soll eine Möglichkeit der Digitalisierung der Umwelt untersucht werden.
In dieser BA wird sich mit einer Lösung zur unabhängigen Umweltdatenerfassung und deren Weiterleitung auf Basis kostengünstiger, im DIY-Bereich verfügbarer, Hardware auseinandergesetzt. Es wird ein PCB - Prototyp gebaut und auf Energiemanagement, Datenerfassung, Konnektivität, Autonomie und Reproduzierbarkeit analysiert.
(Die Ergebnisse Zeigen XYZ)
Die Leiterplatte erfüllt im Fazit die Anforderungen an Datenerfassung, dessen Verarbeitung und Weiterleitung und kann für die OpenSource-Community zur Verfügung gestellt werden.
Die Ergebnisse sollen Maker-Communities, und Kommunen anregen ihre Natürlichen Ressourcen zu Digitalisieren, um eine global verknüpfte Datengrundlage, für den Erhalt der Natur und den Schutz des Planeten zu schaffen.



% \url{https://www.scribbr.de/aufbau-und-gliederung/abstract-schreiben/}



% Fragestellung:
%     Energieautonom
%     Umweltdaten / Sensoren
%     Erfassen
%     Übertragen
% Methoden / Herangehensweise
%     Prototypen Bauen
    
    
% Ergebnisse

% Quellen



    



% Worum geht es?
% Zielsetzung - Menschen ermöglichen ihr Umwelt besser im Blick zu haben.
% Problemstellung - Wie können die Umweltdaten eigenständig erhoben werden, und die Menschen sie einfach einsehen.
% Forschungsfrage - Eine Plattform für die primäre Erfassung, Verarbeitung und Weiterleitung der Umweltdaten. 


% Wie bin ich vorgegangen?
